\section{Introduction}\label{sec:introduction}

A retail-sized payment from a sender in the United States to a recipient in Mexico routed through the SWIFT correspondent banking network carries an all-in cost of approximately 330 basis points of the transferred amount, according to the World Bank's remittance prices panel \citep{rpw_q3_2025}. The same payment routed through a USDT transfer on the Tron network carries a per-transaction fee of approximately 37 basis points at a transfer size of USD 400 \citep{tronsave_blog_2026}. The roughly ninefold gap has been the empirical underpinning of the recent argument that stablecoins represent a frontier technology in cross-border retail payments. The comparison is incomplete, because it prices only the on-chain hop of the stablecoin against the all-in cost of the incumbent.

We organise the paper around a comparison of two payment workflows. On a traditional rail, whether SWIFT correspondent banking or a domestic real-time gross settlement system, the asset the payer briefly holds is a commercial bank deposit: par-redeemable, backstopped by deposit insurance, central bank settlement, and orderly resolution. For a retail-sized payment the payer bears neither counterparty nor depeg risk, so the expected loss on the underlying is zero, and there is no system-level run to price because the public backstops prevent one. A stablecoin payment runs a different workflow: a full fiat-to-fiat circuit in which the sender converts fiat into a private-issuer liability at an on-ramp, holds it however briefly, transfers it on chain, and the recipient converts it back into local fiat at an off-ramp. The rail's operational risk is common to both workflows and nets out of the comparison; what does not net out is the asset the payer holds during the circuit. The expected loss on that liability is the object we price.

The circuit prices leg by leg, and the structure of the cost is the structure of the paper. The on-ramp, the on-chain transfer, and the off-ramp are observable fees. The holding leg carries the expected loss, and it splits into three pieces. The counterparty piece is an issuer-level credit exposure, the product of a default hazard and a loss given default read from market proxies for the issuer's reserve quality; it is observable and requires no model. The baseline depeg piece is the cost of converting the par claim under normal conditions; measured at the secondary market mid it is negligible, and the holder's real cost is the observed off-ramp basis, so it too is read from data. The stress depeg piece is the one element no observable conversion cost reveals: when many holders convert at once, primary redemption is exhausted and the secondary price falls, a self-fulfilling run that has to be priced with a coordination model.

We price that stress piece with a global game in the spirit of \citet{morris_shin_1998}, in the closed-form tradition of \citet{klee_etal_2026_fragility}: uniform signal noise and a linear price impact on a secondary market reached after primary redemption is exhausted at the issuer's liquid reserves. The model delivers analytical objects rather than a high-frequency calibration. The redemption threshold solves a single one-dimensional indifference condition; the expected depeg loss is closed form; the probability of the stress regime is endogenous and specific to each issuer, in the manner of the run probability in \citet{klee_etal_2026_fragility} rather than a frequency shared across coins; and the reserve comparative static, that deeper liquid reserves lower the expected depeg loss, is a theorem rather than a simulation result. The two largest historical episodes calibrate intuition for the two coins: USDC during the March 2023 Silicon Valley Bank failure, when primary redemption was constrained and the price fell to 0.90, and USDT during the May 2022 Terra collapse, when primary redemption stayed open and the price dipped only to 0.975.

The holding leg's expected loss is 21 basis points for USDC and 60 for USDT, both counterparty-led but separated by the severity of the stress depeg, deep for USDC and mild for USDT. The same coordination machinery applied to two issuers with different primary-redemption access returns two distinct risk profiles, so a single supervisory rule would misprice one of them. We then add the full circuit to the on-chain fee and benchmark the result against public, incumbent, and fintech rails across three idealised corridors, a Brazilian domestic payment, a US-Mexico remittance, and an EU-Brazil transfer: the stablecoin is dominated where an upgraded public instant payment rail already exists (Brazil's PIX, at zero cost to the sender), wins where none does, and would in turn be absorbed by an upgraded public cross-border rail, a BIS Project Nexus, whose efficient cost structure would dominate it by one to two orders of magnitude even on the corridors where it currently wins.

Our contribution is to turn the risk exposures of the holding leg into a price, issuer by issuer, with a model that earns its place through mechanism and counterfactuals rather than through the number. The counterparty and baseline pieces are observable; the model prices the coordination run and, because its causal properties are analytical, it answers the policy questions an observable cannot: how the premium moves with reserve quality, secondary market depth, and the regulation that governs them. The two reframed decisions a central bank faces are whether to meet stablecoin adoption by building or by licensing cross-border payment infrastructure, and whether one supervisory rulebook can fit instruments whose risks are composed so differently.

The remainder of the paper is organised as follows. Section~\ref{sec:literature} reviews the related literature. Section~\ref{sec:risk_factors} sets out the holding leg's two risk factors, the counterparty credit exposure and the depeg coordination problem. Section~\ref{sec:model} presents the closed-form global game and its analytical properties. Section~\ref{sec:calibration} parametrises the model by observables and contrasts it with the reduced-form evidence. Section~\ref{sec:empirical_comparison} prices the full circuit across three use cases. Section~\ref{sec:discussion} draws the policy implications. Section~\ref{sec:conclusion} concludes.
